\section{Polynomial factors in the dyadic budget}
\label{sec:area_law_polynomial_decay}

Retain the fixed exponent $\delta_0=2\cdot10^{-7}$. For every real number
$p$ and nonnegative integer $k$, set
\begin{align}
    w_p(k) & =(k+1)^p2^{-\delta_0 k/2}.
\end{align}
The power $(k+1)^p$ is the ordinary real power of a positive number.
The following estimates give the numerical series bound used in
\cite[Section~11]{OpenAI2026AreaLaw}.
To bound the number of repairs by this series, one must also bound the
descendants of each initial mark and relate the initial marks to the
sparse belts.

% Source: 10-geometry.tex, lines 232--235 and 668--692,
% geometry:descendant-count and geometry:total-repairs.
% These are auxiliary numerical statements. The construction and count of
% actual repairs, and the full source statement geometry:total-repairs,
% remain separate obligations.

\begin{theorem}[Eventual absorption of a polynomial factor]
    \label{thm:area_law_polynomial_absorption}
    \lean{TNLean.PEPS.AreaLaw.Geometry.exists_polynomial_dyadic_absorption}
    \leanok
    \uses{thm:area_law_exponent_gaps}
    For every real number $p$, there is a nonnegative integer $K$ such
    that, for every nonnegative integer $k\ge K$,
    \begin{align}
        (k+1)^p2^{-\delta_0 k/2} & \le2^{-\delta_0 k/4}.
    \end{align}
    The threshold depends only on $p$ and the fixed exponent.
\end{theorem}
\begin{proof}\leanok
    \uses{thm:area_law_exponent_gaps}
    Set $a=\delta_0\log2/4>0$. Exponential decay dominates every fixed
    real power, so $(k+1)^p\exp(-ak)$ tends to zero as $k$ tends to
    infinity. Choose $K$ so that this expression is at most one for
    every $k\ge K$. Multiplying by the positive number
    $2^{-\delta_0 k/4}$ gives the asserted inequality.
\end{proof}

\begin{theorem}[Summability of the polynomially weighted decay]
    \label{thm:area_law_polynomial_summability}
    \lean{TNLean.PEPS.AreaLaw.Geometry.summable_polynomial_dyadic_decay}
    \leanok
    \uses{thm:area_law_exponent_gaps}
    For every real number $p$, the nonnegative series
    \begin{align}
        \sum_{k=0}^{\infty}(k+1)^p2^{-\delta_0 k/2}
    \end{align}
    converges.
\end{theorem}
\begin{proof}\leanok
    \uses{thm:area_law_polynomial_absorption,thm:area_law_exponent_gaps}
    Put $q=2^{-\delta_0/4}$. Positivity of $\delta_0$ gives $0<q<1$.
    The absorption estimate bounds $w_p(k)$ by $q^k$ for all sufficiently
    large $k$. The geometric series converges, and adding finitely many
    initial terms preserves convergence.
\end{proof}

\begin{theorem}[Uniform bound on finite scale sums]
    \label{thm:area_law_polynomial_finite_sums}
    \lean{TNLean.PEPS.AreaLaw.Geometry.exists_uniform_polynomial_dyadic_sum_bound}
    \leanok
    \uses{thm:area_law_exponent_gaps}
    For every real number $p$, there is a real number $B>0$ such that,
    for every finite set $F$ of nonnegative integers,
    \begin{align}
        \sum_{k\in F}(k+1)^p2^{-\delta_0 k/2} & \le B.
    \end{align}
    The bound is chosen from $p$ and the fixed exponent alone,
    independently of the origin, domain, cut, and neighborhood radius.
    It therefore applies to every finite selection of scales above any
    chosen lower scale.
\end{theorem}
\begin{proof}\leanok
    \uses{thm:area_law_polynomial_summability}
    Let $S$ be the sum of the convergent nonnegative series and take
    $B=S+1$. Nonnegativity gives $S\ge0$ and bounds every finite subsum
    by $S$, proving both $B>0$ and the stated estimate.
\end{proof}
